\documentclass[12]{article}

\usepackage{import}
\usepackage{transparent}
\usepackage{xcolor}
\usepackage{tcolorbox}
\usepackage{amsmath}
\usepackage{upgreek}
\usepackage{enumitem}
\usepackage{amssymb}
\usepackage[margin=0.25in]{geometry}
\usepackage{listings}
\usepackage[T1]{fontenc}
%\geometry{letterpaper}
%Commands definitions
\newcommand{\setbackgroundcolour}{\pagecolor[rgb]{0.19, 0.19, 0.19}}
\newcommand{\settextcolour}{\color[rgb]{0.87, 0.87, 0.87}}
\newcommand{\invertbackgroundtext}{\setbackgroundcolour\settextcolour}

%Command execution
%If this line is commented, then the appearance remaind as usual.
\invertbackgroundtext

\lstset{
    language=C,
    basicstyle=\ttfamily,
    keywordstyle=\color{cyan},
    commentstyle=\color{green!70!black},
    stringstyle=\color{red},
    showstringspaces=false,
    tabsize=4,
    breaklines=true,
    breakatwhitespace=false,
}

% Ponizsze linijki sa templatka do zadan i problemow                                                                                                           
% Wziete gdzies z internetu
% Assignment class for homework and exams.                                                                                                                  
%
% Gabriel Antonius
%
%
\NeedsTeXFormat{LaTeX2e}
\RequirePackage{placeins}
\RequirePackage{environ}
\RequirePackage{xifthen}


% \solutiontrue and \solutionfalse control whether solutions are shown
\newif\ifsolution
\solutiontrue
%\solutionfalse


% Redefine these if you want to change the language
\newcommand{\problemlabel}{Zadanie}
\newcommand{\solutionlabel}{Rozwi\k{a}zanie}


% Set up counters for problems and subproblems
\newcounter{ProblemNum}
\newcounter{SubProblemNum}[ProblemNum]
\renewcommand{\theProblemNum}{\arabic{ProblemNum}}
\renewcommand{\theSubProblemNum}{\alph{SubProblemNum}}


% The problem environment is the base unit of content for this class.
\newcommand{\subsectiontitle}{}
\newenvironment{problem}[1]%
  {
  \stepcounter{ProblemNum}
  \renewcommand{\subsectiontitle}{\problemlabel \ \theProblemNum \ifthenelse{\isempty{#1}}{}{: #1}}
  \medskip \subsection*{\subsectiontitle}  % new line after title
  {\large Tre\'{s}\'{c}: } 
  \FloatBarrier
  }
  {
  \FloatBarrier
  }


% The subproblem command divides a problem into parts a), b), c), ...
\newcommand*{\subproblem}{%
  \vspace*{3pt}
  \stepcounter{SubProblemNum}%
  {\bf \theSubProblemNum)\hspace{2pt}}
  }


% The solution environment should be used within the problem environment
\NewEnviron{solution}{
  \setcounter{SubProblemNum}{0}
  \ifsolution
    \FloatBarrier
    \subsubsection*{\solutionlabel}
    \BODY
  \fi}


%\documentclass[12pt]{myassignment}
%\usepackage{amsmath}                                                                                                                                        
\setlength{\parindent}{0pt}

\solutiontrue  % <---- Toggle showing solutions or not
%\solutionfalse

%koniec templatki do zadan i problemow




%\author{Lukasz Kopyto}
%\title{tytul}

\begin{document}
%\maketitle
\noindent \L ukasz Kopyto \hfill 05.12.2024 \\

\begin{center}

    {\large Tresc i rozwiazania listy numer 3 z Matematyki Dyskretnej(M)}
\end{center}

\begin{problem}{Hyperbolic functions}

    Niech $f(n) = \sum_{k=1}^{n} \lceil log_{2}k \rceil$

    \subproblem Wykaz, ze $$f(n) = n-1 +f(\lceil\tfrac{n}{2}\rceil) + f(\left\lfloor\tfrac{n}{2}\right\rfloor)$$
dla wszystkich $n \geq 1$


    \subproblem Wykaz, ze jesli w powyzszej zaleznosci wymagamy, aby $f(0) = f(1)$, to f jest jedyna funkcja spelniajaca te zaleznosc.

    {\small (Wskazowka: rozbij sume na n parzyste i nieparzyste)}

    \begin{solution}
\end{solution}
\end{problem}



\begin{problem}{Zwarta postac}
    
    Wyznacz zwarta postac funkcji $f(n)$ z poprzedniego zadania. Zwarta czyli bez symboli $\sum$, \dots, etc
    \begin{solution}
        blebleleb
        
    \end{solution}
\end{problem}

\begin{problem}{Reprezentacja Zeckendorfa}
    Przedstawieniem liczby naturalnej n w ukladzie liczb Fibonacciego, nazywamy taki ciag wspolczynnikow $a_{2}, a_{3}, \dots, a_{k}$ rownych 0 lub 1, ze $a_{2}F_{2} + \dots a_{k}F_{k} = n$ oraz $a_{i} + a_{i+1} \leq 1$ dla wszystkich i.

    Pokaz, ze kazda liczba naturalna ma jednoznaczne przedstawienie w ukladzie liczb Fibonacciego.

    \begin{solution}
        bl
    \end{solution}
    
\end{problem}

\begin{problem}{Potegowanie modularne}

    Niech $x, k, n$ beda liczbami calkowitymi. 
   
    \subproblem Skonstruuj algorytm obliczajacy $x^{k}$modulo n. Algorytm powinien korzystac ze wzorow: $x^{2l} = x^{l}\cdot x^{l}$ oraz $x^{2l+1} = x\cdot x^{2l}$
    
    \subproblem Okresl liczbe mnozen wykonywnaych przez ten algorytm.

    \begin{solution}
        afs
    \end{solution}
\end{problem}

\begin{problem}{Potegowanie macierzy i zwiazane z tym liczby Fibonacciego.}
    
    \subproblem Znajdz wzor na $n$-ta potege macierzy 
    $\begin{bmatrix}
        1 & 1 \\
        1 & 0
    \end{bmatrix}$

    \subproblem Napisz szybki algorytm wyliczajacy z tego wzory $n$-ta liczbe Fibonacciego.

    \subproblem Oszacuj zlozonosc tego algorytmu, jesli uzywa on procedury mnozacej dwie liczby k-cyfrowe w czasie $M(k)$. 

    Zaloz, ze $M(k) \geq 2M(\tfrac{k}{2})$ 

    \begin{solution}
        
    \end{solution}
        
\end{problem}

\begin{problem}{Szybkie potegowanie macierzy oraz uogolnienie wyliczania zaleznosci rekurencyjnej za pomoca macierzy}

    \subproblem Pokaz, ze
    $
    {\begin{bmatrix}
        1 & 1 \\
        1 & 0
    \end{bmatrix}}^{n}
    \begin{bmatrix}
        F_{2} \\
        F_{1}
    \end{bmatrix} = 
    \begin{bmatrix}
        F_{n+2} \\
        F_{n+1}
    \end{bmatrix}
    $
    
    \subproblem Skonstruuj algorytm, ktory dla danych $p_{1},\dots,p_{k},a_{1},\dots,a_{k} \in \mathbb{R}$ wylicza szybko $a_{n}$ zadane zaleznoscia rekurencyjna: $a_{n} = p_{1}a_{n-1} + \dots + p_{k}a_{n-k}$
  
    \begin{solution}
        
    \end{solution}

\end{problem}


\begin{problem}{Mnozenie liczb, algorytm Toom-3?}
    
    \subproblem Skontruuj algorytm mnozacy dwie $n$-cyfrowe liczby calkowite a i b przez rozbicie kazdej z nich na trzy rowne czesci: $a_{1}, a_{2}, a_{3}, b_{1}, b_{2}, b_{3}$. Powinien sie on zadowalac piecioma wywolaniami rekurencyjnymi samego siebie obliczajacymi: 

    $m_{1} = a_{1}\cdot b_{1},$ 

    $m_{2} a_{3} \cdot b_{3}, $

    $m_{3} = (a_{1}+ a_{2}+a_{3})\cdot (b_{1}+b_{2}+b_{3}),$ 

    $m_{4}  = (a_{1}-a_{2}+a_{3})\cdot (b_{1}-b_{2}+b_{3}), $

    $m_{5} = (a_{1} -2a_{2}+4a_{3})\cdot (b_{1}-2b_{2}+4b_{3})$

    \subproblem Oszacuj jego zlozonosc obliczeniowa

    \begin{solution}
        
    \end{solution}

\end{problem}

\begin{problem}{Zlozonosc obliczeniowa $T(n) = a \cdot T(\tfrac{n}{b}) + c\cdot n^{d}$.}
    
    Dany jest algorytm typu "dziel i zwyciezaj" wywolujacy sam siebie rekurencyjnie,$a$ razy dla podproblemow rozmiaru $\frac{n}{b}$ i wykonujacy przy tym  $c\cdots n^{d}$ operacji. Czas $T(n)$ dzialania tego algorytmu spelnia zaleznosc rekurencyjna: $T(n) = a \cdot T(\tfrac{n}{b}) + c\cdot n^{d}$. Korzystajac z tej zaleznosci oszacuj$T(n)$ jako $O(\cdot)$ w zaleznosci od $a,b,d$. Mozesz zalozyc, ze $n = b^{k}$

{\small (Wskazowka: rozwaz trzy przypadki, $a>b^{d}, a = b^{d}, a<b^{d})$}

\begin{solution}
    
\end{solution}

\end{problem}

\begin{problem}{Zlozonosc obliczeniowa  $T(n) \leq T(\lceil \tfrac{n}{5} \rceil + 2) + T(\lceil \tfrac{7n}{10}+ 2\rceil)+ cn.$}
    Niech $T(n) \leq T(\lceil \tfrac{n}{5} \rceil + 2) + T(\lceil \tfrac{7n}{10}+ 2\rceil)+ cn.$

    Pokaz uzywajac indukcji, ze $T(n) \leq c^{\prime}n$ dla pewnej stalej $c^{\prime}$.


    \begin{solution}
        
    \end{solution}
\end{problem}


\begin{problem}{Zaleznosc gcd no.1}
    
    \subproblem Udowonij, ze jesli $a,b \in \mathbb{N}$ oraz x i y sa niezerowymi liczbami calkowitymi spelniajacymi rownanie diofantyczne $ax + by = 1$, to:
    $$gcd(a,b) = gcd(a,y) = gcd(x,b) = gcd(x,y)$$

    \subproblem Wykaz, ze przynajmniej jedna z liczb x i y z powyzszego podpunktu musi byc ujemna.
    
    \begin{solution}
        
    \end{solution}

\end{problem}


\begin{problem}{Ciag dalszy gcd}

    \subproblem Przedstaw $gcd(448,721)$ w postaci $721x + 448y$, dla pewnych $x,y \in \mathbb{Z}$

    \subproblem Oblicz takie calkowite x i y, ze $333x + 1234y = 1$.

    \subproblem Ile sie rowna $333^{-1}$ w pierscieniu $\mathbb{Z}_{1234}$?

    \subproblem Oblicz $-69^{-1}$mod$1313$.
    
    \begin{solution}
        
    \end{solution}

\end{problem}


\begin{problem}{Zaleznosc gcd z liczbami Fibonacciego}
    
    \subproblem Pokaz, ze $gcd(F_{n-1}, F_{n}) = 1$.

    \subproblem Udowodnij indukcyjnie, ze $gcd(F_{m}, F_{n}) = F_{gcd(m,n)}$
    
    
    \begin{solution}
        
    \end{solution}

\end{problem}


\begin{problem}{Liczby wzglednie pierwsze i gcd}

    Udowodnij, ze jesli $a\perp b$ i $a>b$ to $gcd(a^{m}-b^{m}, a^{n}-b^{n})= a^{gcd(m,n)} - b^{gcd(m,n)}, 0 \leq m <n$.
   
    \begin{solution}
        
    \end{solution}

\end{problem}

\begin{problem}{Potega 2 z 1 na czele.}

    Pokaz, ze dla kazdego n istnieje dokladnie jedna potega 2, ktora w ukladzie dziesietnym ma n cyfr z najbardziej znaczaca cyfra 1.
    
    \begin{solution}
        
    \end{solution}


\end{problem}

\begin{problem}{Rownanie chyba diofantyczne}
        
    Niech $ax_{0} + by_{0} = c$ dla pewnych $a,b,c,x_{0}, y_{0} \in \mathbb{Z}$. Okresl zbior wszystkich rozwiazan $(x,y)$ rownania $$ax + by = c$$.


    \begin{solution}
        
    \end{solution}

\end{problem}






\end{document}
